\documentclass[aspectratio=169,11pt]{beamer}
% ANPOF 2026-inspired Beamer theme (unofficial)
% Institutional logo images: unicamp_logo.png and fapesp_logo.jpg.
\usepackage[T1]{fontenc}
\usepackage[utf8]{inputenc}
\usepackage{lmodern}
\usepackage{tikz}
\usetikzlibrary{calc}
\usepackage{amsmath,amssymb,mathtools}
\usepackage{graphicx}
\usepackage{microtype}
\definecolor{ANPOFnavy}{HTML}{204965}
\definecolor{ANPOFdeep}{HTML}{173951}
\definecolor{ANPOForange}{HTML}{D98B46}
\definecolor{ANPOFcream}{HTML}{FFF9E5}
\definecolor{ANPOFgrey}{HTML}{647482}
\setbeamercolor{normal text}{fg=ANPOFdeep,bg=white}
\setbeamercolor{frametitle}{fg=white,bg=ANPOFnavy}
\setbeamercolor{title}{fg=white}
\setbeamercolor{structure}{fg=ANPOForange}
\setbeamercolor{block title}{fg=white,bg=ANPOFnavy}
\setbeamercolor{block body}{fg=ANPOFdeep,bg=ANPOFcream}
\setbeamertemplate{navigation symbols}{}
\setbeamertemplate{itemize item}{\color{ANPOForange}$\blacktriangleright$}
\setbeamersize{text margin left=0.76cm,text margin right=0.76cm}
\setbeamerfont{frametitle}{size=\Large,series=\bfseries}
\newcommand{\shortevent}{XXI ENCONTRO NACIONAL DA ANPOF}

% Institutional logos are included as image files (no TikZ logo sketches).
\newcommand{\InstitutionalLogos}{%
 \begin{minipage}[c]{0.27\linewidth}\centering
 \includegraphics[height=1.13cm,keepaspectratio]{unicamp_logo.png}\end{minipage}%
 \hfill
 \begin{minipage}[c]{0.57\linewidth}\centering
 \includegraphics[width=\linewidth,height=0.92cm,keepaspectratio]{fapesp_logo.jpg}\end{minipage}%
}

% Backgrounds based on the event's dark blue + orange 2026 palette.
\newcommand{\CoverArt}{%
\begin{tikzpicture}[remember picture,overlay]
 \fill[ANPOFnavy] (current page.south west) rectangle (current page.north east);
 \fill[ANPOFdeep] ([xshift=-0.3cm,yshift=0.1cm]current page.north east) circle (2.8cm);
 \fill[ANPOForange] ([xshift=-0.55cm,yshift=-0.65cm]current page.south west) circle (2.6cm);
 \fill[ANPOFnavy] ([xshift=-1.1cm,yshift=-1.1cm]current page.south west) circle (2.6cm);
 \fill[ANPOForange] ([xshift=-0.2cm,yshift=0.1cm]current page.north east) circle (0.23cm);
 \fill[white] (current page.south west) rectangle ([yshift=1.57cm]current page.south east);
 \fill[ANPOForange] ([yshift=1.57cm]current page.south west) rectangle ([yshift=1.62cm]current page.south east);
\end{tikzpicture}}
\setbeamertemplate{frametitle}{%
 \nointerlineskip
 \begin{beamercolorbox}[wd=\paperwidth,ht=1.05cm,dp=.24cm,leftskip=.74cm]{frametitle}
  \insertframetitle
 \end{beamercolorbox}
 \begin{tikzpicture}[remember picture,overlay]
 \fill[ANPOForange] (0,0) rectangle (\paperwidth,0.045cm);
 \end{tikzpicture}
}
\setbeamertemplate{footline}{%
 \ifnum\value{framenumber}>1
 \begin{beamercolorbox}[wd=\paperwidth,ht=.40cm,dp=.12cm,leftskip=.75cm,rightskip=.7cm]{normal text}
 \color{ANPOFgrey}\fontsize{7.8}{9}\selectfont \shortevent\hfill \insertframenumber
 \end{beamercolorbox}
 \fi}

\title{An Axiomatization for Many-Valued\\ Dynamic Epistemic Logics}
\author{Rafael Ongaratto}
\institute{UNICAMP}
\date{Florian\'opolis, 16--20 November 2026}

\begin{document}
\begin{frame}[plain]
\CoverArt
\begin{tikzpicture}[remember picture,overlay]
 \node[anchor=north west,text=ANPOFcream,font=\sffamily\bfseries\fontsize{13}{15}\selectfont,align=left]
 at ([xshift=.82cm,yshift=-.70cm]current page.north west) {XXI ENCONTRO\\NACIONAL DA ANPOF};
 \node[anchor=north east,text=ANPOFcream,font=\sffamily\bfseries\fontsize{9}{11}\selectfont,align=right]
 at ([xshift=-.8cm,yshift=-.79cm]current page.north east) {FLORIAN\'OPOLIS, SC\\16--20 NOV 2026};
 \node[anchor=north west,text=white,text width=11.8cm,font=\sffamily\bfseries\fontsize{25}{29}\selectfont,align=left]
 at ([xshift=.82cm,yshift=-2.20cm]current page.north west) {An Axiomatization for\\Many-Valued Dynamic\\Epistemic Logics};
 \node[anchor=north west,text=ANPOFcream,font=\sffamily\bfseries\fontsize{13}{16}\selectfont]
 at ([xshift=.86cm,yshift=-5.35cm]current page.north west) {Rafael Ongaratto \hspace{0.38cm}\textbar\hspace{0.38cm} UNICAMP};
 \node[anchor=north west,text=white,text width=11.9cm,font=\sffamily\fontsize{9.2}{12}\selectfont,align=left]
 at ([xshift=.86cm,yshift=-5.82cm]current page.north west)
 {\textbf{Supervisor:} Walter Carnielli (CLE, UNICAMP)\quad\\
  \textbf{Co-supervisor:} Hans van Ditmarsch (CNRS/IRIT, Toulouse)};
 \node[anchor=west,inner sep=0pt] at ([xshift=.83cm,yshift=.77cm]current page.south west)
 {\begin{minipage}{6.3cm}\InstitutionalLogos\end{minipage}};
 \node[anchor=east,text=ANPOFdeep,font=\sffamily\fontsize{8.7}{11}\selectfont,align=right]
 at ([xshift=-.7cm,yshift=.82cm]current page.south east)
 {FAPESP Grant No. 2024/00807-6};
\end{tikzpicture}
\end{frame}

\begin{frame}{Outline}
\begin{enumerate}
 \setlength\itemsep{.62cm}
 \item Motivation and background
 \item Many-valued epistemic semantics
 \item Dynamic operators and public announcements
 \item Axiomatization and completeness
 \item Concluding remarks
\end{enumerate}
\end{frame}

\begin{frame}{Motivation and background}
\begin{block}{Many-valued dynamics}
How can dynamic reasoning be axiomatized over many-valued modal logics?
\end{block}
\vspace{0.4cm}
\begin{itemize}
\item Dynamic modalities extend modal reasoning by allowing change over Kripke models. 
\item Classical dynamic epistemic logic is axiomatized by the use of \textit{reduction axioms}, equivalences that reduce every formula of DEL into a formula without dynamic modalities.
\item Completeness of DEL is proved by reduction to the epistemic language. 

\end{itemize}
Problem: can we do the same for many-valued modal logics?
\end{frame}

\begin{frame}{Many-valued epistemic semantics}
\begin{block}{Semantic framework}
% Replace with the precise semantics of the logics used in the talk.
Let $M=(W,R,V)$ be an epistemic model, where $W$ is a set of states,
$R$ is an accessibility relation, and $V$ assigns many-valued valuations.
\end{block}
\vspace{0.5cm}
\centering\color{ANPOFgrey}\emph{Insert truth-value algebra, designated values, and modal clauses here.}
\end{frame}

\begin{frame}{Axiomatization and completeness}
\begin{block}{Main result}
\emph{Insert the precise statement of the axiomatization theorem here.}
\end{block}
\vspace{0.35cm}
\begin{itemize}
 \item Static axioms and inference rules
 \item Dynamic principles / reduction axioms (when applicable)
 \item Soundness and completeness argument
\end{itemize}
\end{frame}

\begin{frame}[plain]
\CoverArt
\begin{tikzpicture}[remember picture,overlay]
 \node[text=white,font=\sffamily\bfseries\fontsize{32}{36}\selectfont]
 at ([yshift=.8cm]current page.center) {Thank you!};
 \node[text=ANPOFcream,font=\sffamily\fontsize{14}{18}\selectfont]
 at ([yshift=-.13cm]current page.center) {Rafael Ongaratto \quad | \quad UNICAMP};
 \node[anchor=west,inner sep=0pt] at ([xshift=.83cm,yshift=.77cm]current page.south west)
 {\begin{minipage}{6.3cm}\InstitutionalLogos\end{minipage}};
 \node[anchor=east,text=ANPOFdeep,font=\sffamily\fontsize{8.7}{11}\selectfont]
 at ([xshift=-.7cm,yshift=.82cm]current page.south east)
 {FAPESP Grant No. 2024/00807-6};
\end{tikzpicture}
\end{frame}
\end{document}
